% =============================================================================
%  3. Estudio Técnico
% =============================================================================
\section{Estudio Técnico}

\subsection{Definición del producto}
(Ver sección~1). Técnicamente, \NombreProducto{} es una aplicación multi-tenant
cloud-native compuesta por un frontend web, una API central, motores de análisis
containerizados y una base de datos de vulnerabilidades.

\subsection{Proceso de producción}

\subsubsection{Subprocesos del servicio}
Al ser software, el «proceso de producción» es el ciclo de desarrollo y operación del
servicio:

\begin{enumerate}
  \item \textbf{Desarrollo del producto:} programación de los módulos (frontend, API,
        motores SAST/SCA/secretos/contenedores, dashboard).
  \item \textbf{Integración del cliente (onboarding):} el cliente conecta su repositorio
        mediante OAuth o instala el agente.
  \item \textbf{Ejecución del análisis:} en cada commit/build, los motores escanean el
        código y generan resultados.
  \item \textbf{Procesamiento e inteligencia:} la capa de IA prioriza vulnerabilidades
        y sugiere remediaciones.
  \item \textbf{Entrega de resultados:} el dashboard muestra hallazgos, alertas y
        reportes de cumplimiento.
  \item \textbf{Soporte y mantenimiento:} atención al cliente, actualización de la base
        de vulnerabilidades y mejoras continuas.
\end{enumerate}

\subsection{Maquinaria y equipo}
En un SaaS no hay maquinaria física de planta; los «equipos» son infraestructura
tecnológica y herramientas:

\begin{table}[ht]
  \centering
  \begin{tabularx}{\linewidth}{l Y}
    \toprule
    \textbf{Componente} & \textbf{Función} \\
    \midrule
    Infraestructura cloud (AWS/Azure) & Servidores, contenedores (Kubernetes), almacenamiento y red \\
    Base de datos gestionada (PostgreSQL) & Almacenamiento de datos de clientes y hallazgos \\
    Motores de análisis (SAST/SCA open-source + propios) & Núcleo del escaneo de seguridad \\
    Feed de vulnerabilidades (CVE, NVD) & Base de conocimiento actualizada \\
    Herramientas de monitoreo y CI/CD propias & Operación y despliegue de la plataforma \\
    Equipos de trabajo del personal (laptops) & Desarrollo y operación \\
    \bottomrule
  \end{tabularx}
\end{table}

\subsection{Tamaño}
\begin{itemize}
  \item \textbf{Capacidad técnica:} arquitectura escalable horizontalmente; soporta
        desde 40 hasta 1\,000+ desarrolladores activos sin rediseño, escalando
        contenedores según demanda.
  \item \textbf{Tamaño objetivo del proyecto (año 5):} 38 empresas cliente,
        $\sim$532 suscripciones activas.
  \item \textbf{Equipo humano inicial:} 3--4 personas (los fundadores cubren
        desarrollo, seguridad, ventas y soporte; se contrata gradualmente).
\end{itemize}

\subsection{Sistemas de control}
\begin{itemize}
  \item \textbf{Control de calidad de software:} pruebas automatizadas, revisión de
        código y la propia plataforma aplicada sobre sí misma (\emph{dogfooding}).
  \item \textbf{Control de seguridad y disponibilidad:} monitoreo 24/7, SLA de
        disponibilidad (99.5\,\%), respaldos y plan de recuperación ante desastres.
  \item \textbf{Control de cumplimiento:} auditorías internas hacia certificación
        ISO 27001 y alineamiento con SOC 2.
  \item \textbf{Control financiero y de uso:} métricas SaaS (MRR, churn, CAC, LTV) y
        medición de consumo por cliente.
\end{itemize}

\subsection{Ingeniería del proceso}
El flujo técnico de valor: \emph{Código del cliente $\rightarrow$ Webhook/Agente
$\rightarrow$ Cola de tareas $\rightarrow$ Motores de análisis en contenedores
$\rightarrow$ Motor de IA de priorización $\rightarrow$ Base de datos $\rightarrow$ API
$\rightarrow$ Dashboard/Alertas/Reportes}. Todo orquestado en Kubernetes con escalado
automático según la carga de escaneos.

\subsection{Plan de producción}
\begin{itemize}
  \item \textbf{Fase 0 (meses 1--4):} MVP con SAST + escaneo de secretos.
  \item \textbf{Fase 1 (meses 5--8):} añadir SCA y dashboard de cumplimiento;
        lanzamiento beta con primeros clientes.
  \item \textbf{Fase 2 (meses 9--12):} escaneo de contenedores/IaC, motor de IA y
        lanzamiento comercial.
  \item \textbf{Año 2 en adelante:} mejora continua, nuevas integraciones y
        escalamiento regional.
\end{itemize}

\subsection{Localización}
\begin{itemize}
  \item \textbf{Sede y operación:} Costa Rica (San José / zona GAM), aprovechando el
        ecosistema tecnológico, talento bilingüe e incentivos de zona franca. [Ref.~8]
  \item \textbf{Infraestructura física del servicio:} en la nube, en regiones de baja
        latencia para Latinoamérica (por ejemplo, regiones de AWS/Azure en EE.\,UU. y
        São Paulo). [Ref.~11]
  \item La localización física de oficinas no es crítica por ser SaaS, pero la sede
        legal en Costa Rica facilita contratación, beneficios fiscales y cercanía con
        el mercado meta.
\end{itemize}

\subsection{Materia prima}
La «materia prima» de un SaaS es intangible:
\begin{itemize}
  \item \textbf{Capital humano} (desarrolladores, ingenieros de seguridad): el insumo principal.
  \item \textbf{Datos de vulnerabilidades} (CVE/NVD, feeds públicos y comerciales).
  \item \textbf{Componentes de software open-source} (motores base como Semgrep, Trivy,
        Gitleaks) integrados y mejorados.
  \item \textbf{Cómputo en la nube} (CPU, memoria, almacenamiento), que se consume por demanda.
\end{itemize}

\subsection{Desarrollo del producto}
El desarrollo es iterativo (metodología ágil/Scrum), con entregas incrementales cada 2
semanas. El roadmap prioriza: (1) cobertura de seguridad, (2) experiencia del
desarrollador, (3) automatización de cumplimiento local y (4) inteligencia artificial
para remediación. El producto se mejora continuamente con base en el feedback de
clientes y la evolución del panorama de amenazas.
